% 型4：段階ヒント型
% 問題のあとにヒントを 3 段階で置き，必要なところまで読んで再挑戦させる．
% 「解答を見てしまう」ことを避けたい層に向く．
\documentclass[lualatex,paper=b5j,fontsize=10pt,twoside]{jlreq}

\usepackage{surikumi-mondaishu}

\MathMagazineSetup{
  feature={},
  series={段階演習／数学 A},
  series-position=left,
  pattern=hexagons,
  title={確率の考え方},
  author={入試基礎講座},
  author-font=mincho,
  author-position=right,
  title-fill=black,
  deck={手が止まったらヒント1だけを読み，もう一度考える．最後まで読まずに解ければそれでよい．}
}

\begin{document}

\MakeMathMagazineTitle

\begin{mmproblem}[1]{}
赤玉 3 個，白玉 5 個が入った袋から，玉を 1 個ずつ取り出す操作を
2 回続けて行う．取り出した玉は袋へ戻さない．
2 回目に赤玉が出る確率を求めよ．
\end{mmproblem}

\MSHint{1}{1 回目に何が出たかで場合が分かれる．2 回目の色だけを見て
数え上げようとすると，袋の中身が定まらない．}

\MSHint{2}{1 回目が赤である事象と白である事象は互いに排反であり，
2 回目に赤が出る事象はこの二つに分けて足し合わせられる．}

\MSHint{3}{1 回目が赤のとき，袋には赤 2 個と白 5 個が残る．
1 回目が白のとき，赤 3 個と白 4 個が残る．}

\MSHintBreak

\begin{mmsolution}
1 回目が赤である事象を $A$，白である事象を $B$ とする．
$A$ と $B$ は互いに排反であり，2 回目に赤が出る事象は
$A$ が起こって赤，または $B$ が起こって赤の場合に分かれる．

$A$ が起こる確率は $\dfrac{3}{8}$ で，このとき袋には赤 2 個と白 5 個が
残るから，続けて赤が出る確率は $\dfrac{2}{7}$ である．

$B$ が起こる確率は $\dfrac{5}{8}$ で，このとき袋には赤 3 個と白 4 個が
残るから，続けて赤が出る確率は $\dfrac{3}{7}$ である．

したがって求める確率は
\begin{align*}
  \dfrac{3}{8}\cdot\dfrac{2}{7}+\dfrac{5}{8}\cdot\dfrac{3}{7}
    &= \dfrac{6}{56}+\dfrac{15}{56}\\
    &= \dfrac{21}{56}\\
    &= \dfrac{3}{8}
\end{align*}
となる．
\end{mmsolution}

\MMNote{答は 1 回目に赤が出る確率と等しい．どの回に取り出しても赤が出る確率は変わらない．}

\MMAlternative{並べ方で数える}
8 個の玉を区別して 1 列に並べる方法は $8!$ 通りある．
このうち 2 番目が赤であるものは，2 番目に置く赤の選び方が 3 通り，
残り 7 個の並べ方が $7!$ 通りであるから $3\cdot 7!$ 通りである．
よって求める確率は
\[
  \dfrac{3\cdot 7!}{8!}=\dfrac{3}{8}
\]
となる．

\end{document}
